% ---------------------------------------------------------------------------
% PDF/UA-oriented build of the linguexx accessibility demo.
%
% On TeX Live 2025-11 or newer, use the first \DocumentMetadata line
% (tagging=on).  On TeX Live 2023-2025 use the second (testphase=phase-III).
% Everything else is identical.  pdfstandard=ua-2 requires PDF 2.0; for a
% PDF 1.7 target use pdfversion=1.7 and pdfstandard=ua-1.
% ---------------------------------------------------------------------------
\DocumentMetadata{lang=en,pdfversion=2.0,pdfstandard=ua-2,tagging=on}
% \DocumentMetadata{lang=en,pdfversion=2.0,pdfstandard=ua-2,testphase={phase-III}}
\documentclass{article}
% T1 under pdfLaTeX: the standard encoding once a document contains
% non-ASCII text, giving single glyphs instead of accents composed onto
% letters.  Recorded so nobody has to re-derive it: on this TeX Live the
% text layer is correct WITHOUT this line too -- the tagging machinery
% supplies the ToUnicode mapping, so copy-paste and screen readers get
% "Mädchen" either way, and veraPDF passes either way.  It is here as
% correct practice for accented input, not as a fix for a demonstrated
% accessibility defect.  The Unicode engines need nothing: they take UTF-8
% natively.
\usepackage{iftex}
\ifPDFTeX \usepackage[T1]{fontenc} \fi
\usepackage[lazy]{linguexx}
\usepackage{hyperref}
\hypersetup{pdftitle={An accessible linguistics document},
            pdfdisplaydoctitle=true}
\GlossTierLang{1}{de}
\title{An accessible linguistics document}
\author{A. Linguist}
\usepackage[normalem]{ulem}

% Not a linguexx concern: the LaTeX kernel's own \footnoterule (the short
% separator above footnotes, article.cls) is pure decoration, but the
% TL2026 tagging=on machinery does not mark it as an Artifact by default,
% so veraPDF flags the \hrule as "neither marked as Artifact nor tagged as
% real content" (PDF/UA-2, 8.2.2).  Reproduces with a bare \documentclass
% {article} and a single \footnote, no linguexx code involved -- an
% upstream kernel gap, worked around here at the document level.
\makeatletter
\ExplSyntaxOn
\renewcommand\footnoterule{%
  \kern-3\p@
  \cs_if_exist:NT \tag_mc_begin:n
    { \tag_if_active:T { \tag_mc_begin:n { artifact } } }
  \hrule\@width.4\columnwidth
  \cs_if_exist:NT \tag_mc_end:
    { \tag_if_active:T { \tag_mc_end: } }
  \kern2.6\p@}
\ExplSyntaxOff
\makeatother

\begin{document}
\maketitle

\section{Tagged linguistic examples}
Ordinary running text precedes the first example.

\ex.\label{intro} A numbered example.
\a. A sub-example.
\b. *A judged sub-example.
\a. A sub-sub-example.
\z.\z.

\exg. Der Hund bellte.\\ the dog bark.\lpzg{pst}\\
\glt `The dog barked.'

% Accented characters in an example, so the T1 line above is exercised and
% not merely present.  With T1 these are single glyphs and the text layer is
% correct; under the OT1 default they would be composed from an accent and a
% letter, look identical on the page, and extract wrongly -- which is what a
% screen reader would then read out.
% Written as literal UTF-8, not as \"a and \ss, because that is what an
% author types -- and under pdfLaTeX it is a different input: inputenc turns
% "ä" into the two tokens \" a, which is how the accents came to collide
% with the sub-example letters in the first place.
\exg. Die Mädchen schliefen draußen.\\
      the girl.\lpzg{pl} sleep.\lpzg{pst} outside\\
\glt `The girls slept outside.'

% Phantom bracket alignment: the gloss word under a bracketed constituent
% lines up with the constituent, not with the bracket.  The \phantom ships
% no ink and no marked content, so the tagging tree is unchanged.
{\GlossPhantomAlign
\exg. ich bin [ein Idiot]\\ I am a idiot\\
\glt `I am an idiot.'}

Examples survive inside footnotes,\footnote{With an example: \ex. Footnoted.

And with a relative reference of its own: \Last\ names that example on the
footnote's own number series.
} and interleave with cross-references such as \ref{intro} --- and with
relative ones such as \Last, which links to the example above it without
needing a label of its own.


% The abbreviations used above, as a tagged list: each label keeps its /E
% expansion, so a screen reader reads "accusative -- accusative".
\section{Abbreviations}
\lpzglist

\section{Glossed alternatives (\texttt{\string\altg}):}

\exg. Die \altg{Frau}{Socke}{Maus}{Tonne} ist hier.\\
      The.\lpzg{sg} \altg{woman.\lpzg{sg}}{sock.\lpzg{sg}}{mouse.\lpzg{sg}}{ton.\lpzg{sg}} is.\lpzg{prs} here.\\

\noindent Two alternatives, multi-word:

\exg. Ich sah \altg{den Mann}{die Frau} gestern.\\
      I see.\lpzg{pst} \altg{the.\lpzg{acc} man}{the.\lpzg{acc} woman} yesterday\\

\noindent\textbf{\texttt{\string\altn} now braces both sides:}

\ex. \altn{\sout{This girl in the red coat}}{She}{Mary} will put a picture of \altn{Bill}{him} on your desk.

\end{document}

%%% Local Variables:
%%% mode: LaTeX
%%% TeX-master: t
%%% End:
